\section{Kết luận}

\subsection{Tổng kết}

Đồ án đã nghiên cứu và xây dựng nguyên mẫu robot trợ lý ảo hỗ trợ chăm sóc, giám sát người cao tuổi tại nhà. Hệ thống kết hợp hội thoại bằng giọng nói, nhắc nhở theo lịch, theo dõi người, phát hiện tình huống nghi té ngã và liên lạc với người chăm sóc. Điện thoại Android gắn trên robot đảm nhiệm thu nhận hình ảnh, âm thanh và xử lý các chức năng hỗ trợ; phần thân robot thực hiện chuyển động theo lệnh từ điện thoại. Cách tổ chức này kế thừa ý tưởng sử dụng điện thoại thông minh làm bộ xử lý của nền tảng OpenBot và được phát triển theo nhu cầu của đề tài.

Hai ứng dụng RobotApp và CaregiverApp tạo thành các phương tiện tương tác dành cho người cao tuổi và người chăm sóc. RobotApp thực hiện hội thoại, phát lời nhắc và xử lý hình ảnh tại robot. CaregiverApp cho phép quản lý thông tin, cấu hình chức năng giám sát, xem nhật ký, tham gia cuộc gọi và điều khiển robot từ xa. Dữ liệu chăm sóc, thông tin cuộc gọi và lệnh chuyển động được trao đổi qua các kênh phù hợp với từng chức năng.

Đối với giám sát té ngã, hệ thống kết hợp ước lượng tư thế, duy trì định danh người mục tiêu và phân tích chuỗi chuyển động bằng LSTM. Kết quả phân loại được kiểm tra bổ sung trước khi robot hỏi xác nhận tình trạng. Khi người dùng cần hỗ trợ hoặc chưa thể xác nhận an toàn, ứng dụng lưu cảnh báo và thực hiện cuộc gọi đến người chăm sóc. Các cơ chế kiểm tra lệnh và dừng khi mất cập nhật được tổ chức ở cả ứng dụng lẫn phần thân robot.

\subsection{Kết quả đạt được}

Những kết quả chính của đồ án gồm:
\begin{itemize}
    \item \textbf{Xây dựng hệ thống hỗ trợ chăm sóc có sự phối hợp giữa robot và người chăm sóc.} Hai ứng dụng cung cấp các chức năng quản lý hồ sơ, nhắc nhở, hội thoại, theo dõi sự kiện và liên lạc từ xa trong cùng một hệ thống.
    \item \textbf{Tích hợp thị giác máy tính với chức năng bám theo và phát hiện té ngã.} Kết quả ước lượng tư thế được sử dụng đồng thời để duy trì người mục tiêu, phân tích chuyển động và điều chỉnh vị trí robot. Bước hậu kiểm và xác nhận bằng giọng nói bổ sung thông tin trước khi phát cảnh báo.
    \item \textbf{Hiện thực khả năng di chuyển và điều khiển từ xa.} Firmware điều khiển độc lập hai động cơ, tiếp nhận lệnh từ RobotApp và đưa vận tốc đặt về 0 khi kết nối hoặc luồng lệnh bị gián đoạn. Việc kiểm tra lệnh được kết hợp giữa ứng dụng và vi điều khiển.
    \item \textbf{Triển khai mô hình AI trên thiết bị Android.} Quy trình chuẩn bị và chuyển đổi đã tạo các mô hình TFLite để tích hợp vào RobotApp. Dữ liệu đầu vào, đầu ra và các phép chuyển đổi tọa độ được tổ chức thống nhất giữa những bước xử lý.
    \item \textbf{Xây dựng các nhóm kiểm thử và bước đầu đánh giá hiệu năng.} Các kiểm thử ở mức logic bao phủ đường dẫn dữ liệu, ánh xạ lệnh động cơ, chính sách bám theo, hình học tư thế và hậu kiểm té ngã. Chương 5 trình bày các chỉ số ban đầu của mô hình LSTM, một cấu hình đo tốc độ suy luận và các kịch bản kiểm thử chức năng đầu cuối.
\end{itemize}

Kết quả hiện tại thể hiện khả năng tích hợp các thành phần của nguyên mẫu. Các số liệu đánh giá ban đầu hỗ trợ xem xét tính khả thi của hướng triển khai, nhưng chưa đại diện cho chất lượng phát hiện té ngã và hiệu năng của toàn bộ hệ thống trong mọi điều kiện sử dụng.

\subsection{Hạn chế}

Phạm vi đánh giá còn giới hạn về dữ liệu, thiết bị và bối cảnh thử nghiệm. Các số liệu đã trình bày chưa đi kèm đầy đủ lịch sử huấn luyện, thông tin chia tập dữ liệu, cấu hình thiết bị và kết quả đo thô để tái lập toàn bộ thí nghiệm. Cấu hình đo tốc độ của biến thể FP16 cũng chưa đại diện cho mọi mô hình được đóng gói trong RobotApp. Vì vậy, cần tiếp tục đánh giá trên đúng phiên bản ứng dụng và thiết bị sử dụng, đặc biệt khi nhiều chức năng hoạt động đồng thời.

Chất lượng giám sát phụ thuộc vào góc nhìn camera, điều kiện chiếu sáng và mức độ che khuất. Dữ liệu té ngã dàn dựng có thể khác với tình huống xảy ra trong sinh hoạt của người cao tuổi. Đặc trưng màu dùng để nhận lại mục tiêu có thể nhầm lẫn khi nhiều người mặc trang phục tương tự; các ngưỡng theo dõi và hậu kiểm cũng cần được hiệu chỉnh theo vị trí camera và tốc độ xử lý. Nguyên mẫu hiện phục vụ hỗ trợ giám sát, chưa được đánh giá lâm sàng và không thay thế thiết bị chẩn đoán y khoa.

Hội thoại qua LLM, đồng bộ dữ liệu từ xa và thiết lập cuộc gọi phụ thuộc vào kết nối Internet. Khi mạng gián đoạn, xử lý hình ảnh và kết nối điều khiển cục bộ không phụ thuộc trực tiếp vào các dịch vụ này, nhưng việc chuyển cảnh báo và liên lạc với người chăm sóc có thể bị gián đoạn. Hai ứng dụng trong một cặp hiện sử dụng cùng tài khoản, nên chưa có cơ chế phân quyền riêng cho robot và người chăm sóc hoặc quản lý nhiều robot trong một gia đình. Cơ chế quản lý khóa truy cập, thông tin xác thực và kiểm tra dữ liệu trên máy chủ cũng cần được hoàn thiện.

Phần thân robot mới đáp ứng các chuyển động cơ bản và cơ chế dừng khi mất lệnh. Việc bổ sung cảm biến tránh va chạm, nút dừng vật lý, giám sát nguồn điện và kiểm tra chặt chẽ chuỗi lệnh nhận vào vẫn cần được xem xét. Dung lượng ứng dụng, mức sử dụng bộ nhớ, nhiệt độ và năng lượng khi vận hành kéo dài cũng chưa được tối ưu đầy đủ.

\subsection{Hướng phát triển}

\subsubsection{Hoàn thiện dữ liệu và phương pháp đánh giá}

Hướng ưu tiên là chuẩn hóa quy trình thí nghiệm để có thể lặp lại và so sánh giữa các phiên bản. Dữ liệu cần được phân chia theo video và người tham gia, đồng thời lưu cấu hình huấn luyện, phiên bản mô hình và kết quả đo. Tập kiểm tra nên mở rộng các tình huống sinh hoạt dễ nhầm với té ngã, nhiều góc camera và điều kiện che khuất. Việc đánh giá cần bao gồm tỷ lệ phát hiện đúng, cảnh báo nhầm và thời gian phản ứng ở cấp sự kiện. Các thí nghiệm lần lượt loại bỏ từng thành phần, chẳng hạn bộ hậu kiểm hoặc thông tin tư thế, sẽ giúp xác định đóng góp của chúng đối với kết quả chung.

\subsubsection{Nâng cao độ ổn định của theo dõi và phát hiện té ngã}

Chức năng theo dõi có thể được cải thiện bằng đặc trưng ngoại hình phù hợp với tài nguyên của điện thoại, kết hợp vị trí và diễn biến chuyển động để nhận lại người mục tiêu. Việc sử dụng độ tin cậy của các điểm khớp và khoảng thời gian thực giữa các khung hình cũng cần được nghiên cứu, nhất là khi tốc độ xử lý thay đổi. Ngoài camera, có thể khảo sát thêm cảm biến khoảng cách hoặc thiết bị đeo để bổ sung thông tin trong những tình huống quan sát bị hạn chế. Mỗi phương án cần được đánh giá về hiệu quả, chi phí tính toán và mức độ thuận tiện đối với người cao tuổi.

\subsubsection{Hoàn thiện khả năng di chuyển và an toàn vận hành}

Phần điều khiển cần được mở rộng theo hướng kiểm tra đầy đủ dữ liệu nhận vào, giới hạn tốc độ thay đổi vận tốc và quy định rõ mức ưu tiên giữa dừng khẩn cấp, điều khiển từ xa và bám theo tự động. Cảm biến khoảng cách, cơ cấu phát hiện va chạm, giám sát pin và nút dừng vật lý là những thành phần có thể bổ sung để hỗ trợ vận hành trong nhà. Các thay đổi cần được đánh giá qua thử nghiệm cơ khí, thời gian dừng và phản ứng khi kết nối bị gián đoạn.

\subsubsection{Mở rộng quản lý người dùng và khả năng duy trì liên lạc}

Cấu trúc quản lý có thể mở rộng để mỗi người chăm sóc sử dụng tài khoản riêng và được cấp quyền truy cập vào robot phù hợp. Thông tin cuộc gọi, lịch sử và lệnh điều khiển cần gắn với từng phiên hoặc thiết bị để tránh nhầm lẫn khi số lượng người dùng tăng. Bên cạnh đó, cơ chế xác thực dịch vụ và bảo vệ dữ liệu cần được hoàn thiện.

Để giảm ảnh hưởng của mất mạng, có thể lưu tạm cảnh báo tại thiết bị và đồng bộ lại khi kết nối được khôi phục, kèm cơ chế tránh tạo bản ghi trùng. Các câu hỏi xác nhận mẫu và phản hồi bằng giọng nói tại thiết bị cũng cần được chuẩn bị cho trường hợp dịch vụ trực tuyến không khả dụng. Khả năng sử dụng cuộc gọi hoặc tin nhắn qua mạng di động làm kênh dự phòng có thể được khảo sát trên thiết bị hỗ trợ.

\subsubsection{Cải thiện trải nghiệm và tối ưu tài nguyên}

Giao diện và tương tác giọng nói cần được đánh giá với người dùng mục tiêu, tập trung vào kích thước chữ, độ tương phản, tốc độ đọc, cách xưng hô, phương ngữ và thời gian chờ phản hồi. Robot cần thông báo rõ trạng thái sử dụng camera, microphone và cho phép người dùng chủ động bật hoặc tắt các chức năng.

Về hiệu năng, cần so sánh các phương án suy luận trên CPU và bộ tăng tốc ở đúng thiết bị sử dụng, đồng thời theo dõi bộ nhớ, nhiệt độ và mức tiêu thụ năng lượng khi chạy lâu. Việc giảm kích thước mô hình, tái sử dụng bộ đệm ảnh và loại bỏ thư viện không cần thiết có thể giúp giảm tải cho điện thoại. Mỗi thay đổi tối ưu cần được kiểm tra lại về chất lượng nhận dạng và độ ổn định của ứng dụng.

\subsection{Kết luận cuối cùng}

Đồ án tạo cơ sở để tiếp tục nghiên cứu việc kết hợp tương tác giọng nói, giám sát bằng thị giác và robot di động trong hỗ trợ chăm sóc tại nhà. Giá trị của nguyên mẫu nằm ở khả năng phối hợp các chức năng này để người cao tuổi nhận được hỗ trợ hằng ngày và người chăm sóc có thêm phương tiện theo dõi, liên lạc từ xa. Những kết quả và giới hạn đã xác định giúp định hướng các bước thử nghiệm, cải tiến tiếp theo, hướng đến giải pháp phù hợp hơn với nhu cầu sử dụng của người cao tuổi.
